\documentclass[10pt,a4paper]{amsart}
\usepackage[margin=19mm]{geometry}
\usepackage{amsmath,amssymb,mathtools}
\usepackage[scr=rsfs,cal=euler]{mathalfa}
\usepackage{xfrac}
\usepackage[unicode,hidelinks,bookmarks=false]{hyperref}
\newcommand{\ie}{i.e.\ }
\newcommand{\cf}{cf.\ }
\newcommand{\resp}{respectively\ }
\newcommand{\Subsection}{\subsection}
\providecommand{\nobreakdash}{\nobreak\hbox{-}}
\setlength{\emergencystretch}{2em}
\multlinegap0pt
\usepackage{bm}
\usepackage{enumerate}
\usepackage{amssymb}
\usepackage{float}
\newtheorem{lem}{Lemma}
\newtheorem{thm}{Theorem}
\title{Proof of a conjecture on graph polytope}
\author{Feihu Liu}
\date{Source: arXiv:2409.11970v2 (CC BY 4.0)}
\begin{document}
\maketitle

\noindent\emph{Source and licence.} Proof of a conjecture on graph polytope. Version arXiv:2409.11970v2 (CC BY 4.0), distributed by arXiv under the Creative Commons Attribution 4.0 International licence, http://creativecommons.org/licenses/by/4.0/, as recorded in the arXiv licence metadata for this exact version and shown on the abstract page https://arxiv.org/abs/2409.11970v2. Full licence notice: \texttt{licenses/arxiv-2409.11970-CC-BY-4.0.txt}.\par\smallskip

\noindent\textit{Editorial scope.} Counted unit: the article's thm UniformHyperGr (source0.tex lines 211--213). The article's complete original proof of it is delivered with it (source0.tex lines 343--383, one \texttt{proof} environment of 41 lines, which the article places away from the statement). No mathematics is written, repaired or re-derived: the statement and the proof are the article's own lines, in the article's own notation, and no retained line is substituted or re-flowed.  Retained uncounted as the source-local context the proof needs: lines 173--181; lines 234--239.  Nothing was omitted from any retained span, so the package carries no omission record.  The retained lines cite 2 works, whose entries are transcribed verbatim from the article's own bibliography source in the e-print and delivered with short editorial labels recorded in the item's reference mapping.  No retained line contains a figure environment, an image-inclusion command or drawing code, so no image is delivered and the package declares no asset dependency.  Editorial formatting only, in addition: an AMS \texttt{amsart} preamble that restates the article's own declarations and macros used by the retained lines, namely the environments \texttt{lem}, \texttt{thm} on separate counters, together with the standard packages those lines need.  The retained environments print in this document as Lemma 1, Theorem 1 in order of appearance, whereas the article prints the same items with the numbering of its own full document.  The complete native source of the article as distributed in the arXiv e-print for this exact version is bundled unchanged as \texttt{sources/165451/source0.tex}.
\begin{lem}\cite{Ehrhart62}\label{Ehrhartpoled+1}
Let $\mathcal{P}$ be a rational polytope of dimension $d$ in $\mathbb{R}^m$ with vertex set $\widehat{V}$. The Ehrhart series $\mathrm{Ehr}(\mathcal{P},x)$ is a rational function $\frac{R(x)}{T(x)}$ where $R(x)$ and $T(x)$ are polynomials satisfying:
\begin{enumerate}
    \item $T(x) = \prod_{\gamma \in \widehat{V}} (1 - x^{\text{den}\gamma})$ for some representation,
    \item When written in lowest terms, $\deg(R(x)) < \deg(T(x))$,
    \item $x=1$ is a pole of order exactly $d+1$,
    \item No pole has order exceeding $d+1$.
\end{enumerate}
\end{lem}

\begin{thm}[Stanley's reciprocity theorem, \cite{StanleyMagic73}]\label{StanleyReciprocity}
Let $\Phi$ be an $r \times m$ integer matrix of full rank $r$. If there exists a vector $\boldsymbol{\alpha} \in \mathbf{nul}(\Phi) \cap \mathbb{P}^m$, then $F(\mathbf{x})$ and $\overline{F}(\mathbf{x})$ are rational functions in the variables $x_1, \dots, x_m$ satisfying
\[
\overline{F}(x_1, x_2, \ldots, x_m) = (-1)^{m-r} F(x_1^{-1}, x_2^{-1}, \ldots, x_m^{-1}).
\]
\end{thm}

\begin{thm}\label{UniformHyperGr}
Let $H=(V,E^*)$ be a finite simple connected hypergraph with vertex set $V=\{v_1,v_2,\ldots, v_k\}$ and hyperedge sequence $E^*=(e_1,e_2,\ldots,e_r)$ such that $\cup_{i=1}^r e_i=V$. If $H$ is $s$-uniform, then $\mathrm{Ehr}(P(H),x)$ is a rational function that, when written in lowest terms, takes the form $\frac{M(x)}{Q(x)}$, where $Q(x)$ is a polynomial and $M(x)$ is a symmetric polynomial of degree $\deg Q(x) - (s+1)$.
\end{thm}

\begin{proof}[Proof of Theorem~\ref{UniformHyperGr}]
By the definition of the hypergraph polytope, we obtain the system:
\begin{align*}
&n_{i_1} + n_{i_2} + \cdots + n_{i_s} \leq t, \quad \text{if  } e_\ell=\{v_{i_1}, v_{i_2}, \ldots, v_{i_s}\} \in E,\quad 1 \leq \ell \leq r, \\
\Longleftrightarrow \quad & n_{i_1} + n_{i_2} + \cdots + n_{i_s} + b_{\ell} - t = 0,\ \text{where } b_{\ell} \geq 0,\ 1 \leq \ell \leq r,\ \text{if }e_\ell=\{v_{i_1}, v_{i_2}, \ldots, v_{i_s}\} \in E.
\end{align*}
Denote this homogeneous linear Diophantine system by $\Phi \mathbf{x} = \mathbf{0}$, where $\Phi$ is an $r \times (k + r + 1)$ matrix of full rank $r$. The generating functions $F(\mathbf{x})$ and $\overline{F}(\mathbf{x})$ are defined as
\[
F(\mathbf{x}) = \sum_{(n_1,\ldots,n_k, b_1,\ldots,b_r,t) \in \mathbf{nul}(\Phi) \cap \; \mathbb{N}^{k+r+1}} y_1^{n_1} \cdots y_k^{n_k} z_1^{b_1} \cdots z_r^{b_r} x^t
\]
and
\[
\overline{F}(\mathbf{x}) = \sum_{(n_1,\ldots,n_k, b_1,\ldots,b_r,t) \in \mathbf{nul}(\Phi) \cap \; \mathbb{P}^{k+r+1}} y_1^{n_1} \cdots y_k^{n_k} z_1^{b_1} \cdots z_r^{b_r} x^t.
\]

Note that $\mathrm{Ehr}(P(H), x) = F(\mathbf{1}, x)$. Observe that the vector $\boldsymbol{\beta} = (\boldsymbol{1}, \boldsymbol{1}, s+1)$ lies in $\mathbf{nul}(\Phi) \cap \mathbb{P}^{k+r+1}$ (where the first $\boldsymbol{1}$ has length $k$, the second length $r$). Thus,
\[
\boldsymbol{\alpha} \in \mathbf{nul}(\Phi) \cap \mathbb{P}^{k+r+1} \quad \Longleftrightarrow \quad \boldsymbol{\alpha} - \boldsymbol{\beta} \in \mathbf{nul}(\Phi) \cap \mathbb{N}^{k+r+1}.
\]
Consequently,
\[
\overline{F}(\mathbf{x}) = y_1 \cdots y_k z_1 \cdots z_r x^{s+1} F(\mathbf{x}).
\]
By Stanley's reciprocity theorem (Theorem~\ref{StanleyReciprocity}),
\[
x^{s+1} F(\mathbf{1}, x) = \overline{F}(\mathbf{1}, x) = (-1)^{k+1} F(\mathbf{1}, x^{-1}),
\]
that is,
\begin{align}\label{EHRPGMH}
x^{s+1} \mathrm{Ehr}(P(H), x) = (-1)^{k+1} \mathrm{Ehr}(P(H), x^{-1}).
\end{align}


Equation~\eqref{EHRPGMH} is therefore equivalent to
\[
M(x) \cdot x^{s+1} = (-1)^{k+1} M(x^{-1}) \frac{Q(x)}{Q(x^{-1})}.
\]
By Lemma~\ref{Ehrhartpoled+1} and $\dim P(H)=k$, we have $M(x) = x^{\deg Q(x) - s - 1} M(x^{-1})$.
Thus, the numerator $M(x)$ is a symmetric polynomial of degree $\deg Q(x) - s - 1$.
This completes the proof.
\end{proof}

\begin{thebibliography}{99}
\bibitem[Ehrhar]{Ehrhart62} E. Ehrhart, \emph{Sur les polyh\'edres rationnels homoth\'etiques \'a $n$ dimensions}, C. R. Acad Sci. Paris. 254 (1962), 616--618.
\bibitem[Stanle]{StanleyMagic73} R. P. Stanley, \emph{Linear homogeneous diophantine equations and magic labelings of graphs}, Duke Math. J. 40 (1973), 607--632.
\end{thebibliography}
\end{document}
